% !TEX program = LuaLaTeX+se
% !TEX encoding = UTF-8 Unicode

\documentclass[12pt, a4paper]{article}
\usepackage[latin, polish]{babel} % ustawienia języka
\usepackage[T1]{fontenc} % przełączenie na fonty EC (256 znaków)
\usepackage[utf8]{luainputenc} % kodowanie
\usepackage{fullpage} % redukcja marginesów
\usepackage{setspace} % odstępy między liniami
\usepackage{geometry} 
\geometry{
	left=12mm,
	right=12mm,
	top=12mm,
	bottom=12mm,
	marginparwidth=0mm,
	marginparsep=0mm
	}

\usepackage{fontspec}
\usepackage{libertine} % ust. czcionki
\pagestyle{empty}	% usuwanie nagłówków/stopek z dokumentu
\usepackage[indent=0pt]{parskip} % ustawianie paragrafów
\setstretch{1.5} 

%  Konfiguracja kolumn
\usepackage{paracol}
\setcolumnwidth{8.3cm,4.3cm}
\setlength{\columnsep}{0.4cm}

\usepackage{gregoriotex}
\gresetlinecolor{gregoriocolor}
\grechangestyle{initial}{\setmainfont{VectisMonumentalW01-Regular}\fontsize{23}{23}\selectfont}
% \grechangedim{spacelinestext}{0.7 cm}{scalable}
% \grechangedim{spaceabovelines}{0.43 cm}{scalable}
% \gresetheadercapture{commentary}{grecommentary}{string}
% \grecommentary{}
\gresetheadercapture{annotation}{greannotation}{string}
% Format annotation above initial
\grechangestyle{annotation}{\small\bfseries}%
%\grechangedim{intersyllablespacenotes}{0.18cm}{scalable}%
%\grechangestyle{commentary}{\it}
%\gresetmodenumbersystem{roman-majuscule}
\usepackage{titlesec} % ustawianie nagłówków
\titlespacing*{\section}{0pt}{0pt}{3mm}	
\titleformat{\section}[hang]{\filcenter\large\scshape\bfseries}{}{0pt}{} % format nagłówków

\newfontface\vectis{VectisMonumentalW01-Regular}
\def\greinitialformat#1{%
{\fontsize{24}{24}\vectis #1}}

\titlespacing*{\section}
{0pt}{0pt}{0pt}

\begin{document}
\large
\begin{spacing}{0.7}
\begin{center}
	\section*{Te saeculorum Principem}
	\small\textit{Liber antiphonarius, Solesmes, 1960, p. 894.2}
	\vspace{15mm}
\end{center}	
\end{spacing}
\greannotation{Hymn.} 
\greannotation{1.} 
\gabcsnippet{(c4) Te(cd) sae(ixdh'!iv)cu(h')ló(h)rum(hg) Prín(gh)ci(g)pem,(g_f) (,) Te,(h) Chri(h')ste,(g) Re(fe)gem(d) Gén(ef)ti(dc)um,(cd..) (;) Te(a) mén(cd)ti(d)um,(df) Te(fd) cór(e')di(d)um(c_') (,) U(g)num(g') fa(f)té(ghg)mur(ed) ár(ef)bi(dc)trum.(cd..) 2.(::z) Sce(cd)lé(ixdh'!iv)sta(h') tur(h)ba(hg) clá(gh)mi(g)tat:(g_f) (,) Re(h)gná(h')re(g) Chri(fe)stum(d) nó(ef)lu(dc)mus:(cd..) (;) Te(a) nos(cd) o(d)ván(df)tes(fd) ó(e')mni(d)um(c_') (,) Re(g)gem(g') su(f)pré(ghg)mum(ed) dí(ef)ci(dc)mus.(cd..) 3.(::z) O(cd) Chri(ixdh'!iv)ste,(h') Prin(h)ceps(hg) Pá(gh)ci(g)fer,(g_f) (,) Men(h)tes(h') re(g)bél(fe)les(d) súb(ef)ji(dc)ce:(cd..) (;) Tu(a)ó(cd)qu<i>e</i>() a(d)mó(df)re(fd) dé(e')vi(d)os,(c_') (,) O(g)vi(g')l<i>e</i>() in(f) u(ghg)num(ed) cón(ef)gre(dc)ga.(cd..) 4.(::z) Ad(cd) hoc(ixdh'!iv) cru(h')én(h)t<i>a</i>() ab(hg) á(gh)bo(g)re,(g_f) (,) Pen(h)des(h') a(g)pér(fe)tis(d) brá(ef)chi(dc)is:(cd..) (;) Di(a)rá(cd)que(d) fos(df)sum(fd) cú(e')spi(d)de(c_') (,) Cor(g) i(g')gne(f) fla(ghg)grans(ed) éx(ef)hi(dc)bes.(cd..) 5.(::z) Ad(cd) hoc(ixdh'!iv) in(h') a(h)ris(hg) áb(gh)de(g)ris(g_f) (,) Vi(h)ni(h') da(g)pí(fe)squ<i>e</i>() i(d)má(ef)gi(dc)ne,(cd..) (;) Fun(a)dens(cd) sa(d)lú(df)tem(fd) fí(e')li(d)is(c_') (,) Trans(g)ver(g')be(f)rá(ghg)to(ed) pé(ef)cto(dc)re.(cd..) 6.(::z) Te(cd) na(ixdh'!iv)ti(h')ó(h)num(hg) Praé(gh)si(g)des(g_f) (,) Ho(h)nó(h')re(g) tol(fe)lant(d) pú(ef)bli(dc)co,(cd..) (;) Co(a)lant(cd) ma(d)gí(df)stri,(fd) jú(e')di(d)ces,(c_') (,) Le(g)ges(g') et(f) ar(ghg)tes(ed) éx(ef)pri(dc)mant.(cd..) 7.(::z) Sub(cd)mís(ixdh'!iv)sa(h') re(h)gum(hg) fúl(gh)ge(g)ant(g_f) (,) Ti(h)bi(h') di(g)cá(fe)t<i>a</i>() in(d)sí(ef)gni(dc)a:(cd..) (;) Mi(a)tí(cd)que(d) sce(df)ptro(fd) pá(e')tri(d)am(c_') (,) Do(g)mó(g')sque(f) sub(ghg)de(ed) cí(ef)vi(dc)um.(cd..) 8.(::z) Je(cd)su,(ixdh'!iv) ti(h')bi(h) sit(hg) gló(gh)ri(g)a,(g_f) (,) Qui(h) sce(h')ptra(g) mun(fe)di(d) tém(ef)pe(dc)ras,(cd..) (;) Cum(a) Pa(cd)tr<i>e</i>,() et(d) al(df)mo(fd) Spí(e')ri(d)tu,(c_') (,) In(g) sem(g')pi(f)tér(ghg)na(ed) saé(ef)cu(dc)la.(cd..) (::z) A(ded)men.(cd..) (::z)}

\vspace{5mm}
\setstretch{1}% odstęp między liniami
1. Ufamy Tobie, o Chryste, głosimy, żeś władcą wieków, Królem narodów i Panem, zarządcą serc i umysłów.\\
2. Dyszący złością tłum woła: „Nie chcemy, by nam był królem”, a my Ciebie królem zwiemy, co wszystko w ręku swym trzyma.\\
3. Jezu, nasz Książe pokoju, panuj nad duchem złowrogim, miłości głosem sprowadzaj owce do jednej owczarni.\\
4. Dla nich ramiona rozpiąłeś, na drzewie krwią pomazanym, a serce ostrzem otwarte spłonęło ogniem miłości.\\
5. To dla nich w chleba postaci i w winie skryłeś oblicze, aby zbawieniem nakarmić z serca otwartej rany.\\
6. Niech Ciebie władcy narodów na placach pieniem swym chwalą, uczeni, sędziowie niech wielbią, niech prawa i prace sławią.\\
7. Niechaj insygnia królewskie błyszczą pod Twymi stopami, niech miłe Twoje władanie obejmie ojczyznę całą.\\
8. Niech chwała, Jezu, Ci będzie, Co świata rządzisz berłami, Wraz z ojcem, Duchem Najświętszym Na wieczne wieki bez końca.\footnote{tłum. Michał Gołębiowski}
\end{document}